% ============================================================================
\clearpage
\section*{v4.9.0-r02 Source-Only: The Round Trip about the Photon and the Umkehrwalze}
\addcontentsline{toc}{section}{v4.9.0-r02 Source-Only: The Round Trip about the Photon and the Umkehrwalze}
\label{sec:v490-r02-round-trip-umkehrwalze}
% ============================================================================

\paragraph{Purpose.}
Theoremlet r19.2 records a \(\mathbb Z_2\) holonomy of the duality circle around the photon.  This layer follows a round
trip matter \(\to\) photon \(\to\) tachyon \(\to\) photon \(\to\) matter in Reading~B, asks on which side of the photon each
passage runs, and asks what the Umkehrwalze of the main body is in physical terms.  It uses the reflector and leakage
language of the OP5 layer (Definition~\ref{def:working-shared-matrix-leakage}).  The computations are
\texttt{llm/k38\_umkehrwalze.py}.  No value is compared with data.  Earlier layers are not rewritten.

\noindent\textbf{Source-only status.}
This v4.9.0-r02 layer is source-only.  It does not build a PDF, does not assign a new DOI,
and does not authorize a public release.  The v4.9.0 Version DOI
\href{https://doi.org/10.5281/zenodo.23534954}{10.5281/zenodo.23534954} and the Concept DOI
\href{https://doi.org/10.5281/zenodo.19210728}{10.5281/zenodo.19210728} are unchanged by it.

\paragraph{Author's motivation (2026-10-21).}
The Umkehrwalze was introduced for a path that does not end at the tachyon but returns, matter \(\to\) photon \(\to\)
tachyon \(\to\) photon \(\to\) matter, read as a parallel transport over the sphere.

\subsection*{v4.9.0-r02 the two sides of the photon}
\label{subsec:v490-r02-sides}

Throughout, \(z=E+iB\), \(s=z\cdot z=E^2-B^2+2i\,E\cdot B\), matter \(s>0\), tachyon \(s<0\), photon \(s=0\)
(\S\ref{subsec:v480-r19-transition-group}).

\paragraph{Theoremlet v4.9.0-r02.1.}
A passage from matter to tachyon at fixed \(|s|\) runs through the upper or the lower half of the \(s\)-plane and carries
\(z\mapsto iz\) or \(z\mapsto -iz\).  The side is the sign of \(\operatorname{Im}s=2\,E\cdot B\); it is Lorentz-invariant,
odd under \(P\) and \(T\), and even under \(C\).  For the round trip:
\begin{enumerate}[leftmargin=2em]
\item[(a)] out and back on the same side (the return retraces the outward path): \(z\mapsto z\);
\item[(b)] out on one side and back on the other (the duality angle keeps its sense, \(\pi/2+\pi/2\)): \(z\mapsto -z\).
\end{enumerate}
On the spinor of Theoremlet r06.4 the round trip (b) is \(\psi\mapsto\pm i\psi\); two round trips give
\(\psi\mapsto-\psi\) with \(z\) restored.

\paragraph{Proof sketch.}
\(s(e^{i\alpha}z)=e^{2i\alpha}s\), so \(\alpha\in[0,\pi/2]\) moves \(s\) through the upper half-plane and
\(\alpha\in[-\pi/2,0]\) through the lower one.  \(SO(3,\mathbb C)\) preserves \(s\) and hence \(\operatorname{Im}s\).  By the
parities of \(E\) and \(B\) (\(P\): \(E\) odd, \(B\) even; \(T\): \(E\) even, \(B\) odd; \(C\): both odd), \(E\cdot B\) is
\(P\)- and \(T\)-odd and \(C\)-even.  Case (b) is a loop around \(s=0\), case (a) is not (Theoremlet r19.2).  The spinor
statement is \(z(e^{i\varphi}\psi)=e^{2i\varphi}z(\psi)\).  Checked numerically.  \(\square\)

\paragraph{Reading.}
The ``double tunnel point'' of the round trip is one point with two sides: the photon, passed on the side
\(E\cdot B>0\) or \(E\cdot B<0\).  The side is the sign of the density of the electromagnetic \(\theta\)-term
\(F\tilde F\).  Only a round trip that changes side encloses the photon.

\subsection*{v4.9.0-r02 the round-trip Umkehrwalze is charge conjugation}
\label{subsec:v490-r02-charge-conjugation}

\paragraph{Theoremlet v4.9.0-r02.2.}
On the field strength the holonomy \(z\mapsto-z\) of the round trip (b) is charge conjugation \(C\), equals \(PT\)
(with \(x\mapsto-x\)), and equals the duality rotation by \(\pi\).  It commutes with \(SO(3,\mathbb C)\) and with the duality
circle, and fixes no non-zero field.

\paragraph{Proof sketch.}
\(C\): \(E\mapsto-E\), \(B\mapsto-B\).  \(T\): \(z\mapsto\bar z\); \(P\): \(z\mapsto-\bar z\); composed, \(z\mapsto-z\).
\(e^{i\pi}=-1\) is central.  \(\square\)

\paragraph{Consequence.}
The Umkehrwalze of the round trip is not an added element.  It is forced by the two-to-one relation between \(z\) and
\(s\), whose branch point is the photon, and in physical terms it is \(C\).  Being central, it cannot serve as the \(J\) of
the readout projector \(C_{\HC}=\tfrac12(\mathrm{id}+J)\) of Definition~\ref{def:working-shared-matrix-leakage}: for
\(J=-1\) that projector is zero.

\subsection*{v4.9.0-r02 the involutive Umkehrwalze must contain time reversal}
\label{subsec:v490-r02-antilinear}

The main body uses \(J\) as an involution that exchanges matter and tachyon (Theorem~\ref{thm:core}, where the inside
view runs with time reversed).  In Reading~B the quarter turn \(z\mapsto iz\) does this exchange, but its square is
\(-1\), so it is not an involution.

\paragraph{Theoremlet v4.9.0-r02.3.}
\begin{enumerate}[leftmargin=2em]
\item[(a)] No \(\mathbb C\)-linear involution \(L\) of \(\mathbb C^3\) satisfies \(s(Lz)=-s(z)\).
\item[(b)] The \(\mathbb C\)-antilinear involutions with \(s(Lz)=-\overline{s(z)}\) are exactly \(Lz=iR\bar z\) with
\(R\in O(3,\mathbb C)\) and \(R\bar R=1\).  Instances are \(R=1\), real orthogonal involutions, and squares of boosts.
\item[(c)] The instance \(R=1\),
\[
J_B z:=i\bar z,\qquad\text{that is}\quad E\mapsto B,\ B\mapsto E,
\]
is the duality rotation by \(\pi/2\) after time reversal.  It exchanges matter and tachyon, maps the null cone to
itself, and preserves the side \(\operatorname{sign}\operatorname{Im}s\).
\end{enumerate}

\paragraph{Proof sketch.}
(a) \(M:=-iL\) preserves \(s\), so \(M\in O(3,\mathbb C)\); \(L^2=1\) gives \(M^2=-1\) and \(\det(M)^2=-1\), impossible.
(b) \(Lz=A\bar z\) with \(A^{\mathsf T}A=-1\) gives \(A=iR\), \(R\in O(3,\mathbb C)\); \(L^2=A\bar A=R\bar R\).
(c) \(i(E-iB)=B+iE\), and \(s\mapsto-\bar s\).  Checked numerically.  \(\square\)

\paragraph{Consequence.}
In three complex dimensions an involutive exchange of matter and tachyon cannot be linear.  It has to be antilinear,
and antilinearity is time reversal.  This matches the inside view of Theorem~\ref{thm:core}, which runs with time
reversed.  The physical reading of the involutive Umkehrwalze is therefore electric--magnetic exchange, a duality
quarter turn composed with \(T\).  It is a symmetry of the source-free Maxwell equations.  Unlike \(C\) it is not
Lorentz-invariant: \(J_B\Lambda=\bar\Lambda J_B\), so it commutes with real rotations only and singles out a frame.

\paragraph{Theoremlet v4.9.0-r02.4.}
The eigen-channels of \(J_B\) lie one on each side of the photon:
\[
J_Bz=z\iff E=B,\ s=2i|E|^2;\qquad J_Bz=-z\iff E=-B,\ s=-2i|E|^2.
\]
The projector \(C_{\HC}=\tfrac12(\mathrm{id}+J_B)\) maps every field to \(\operatorname{Re}s=0\),
\(\operatorname{Im}s\ge0\).  Its image meets the null cone only in \(z=0\).

\paragraph{Proof sketch.}
Direct substitution; checked on 1000 random fields (\(\max|\operatorname{Re}s|=0\), \(\min\operatorname{Im}s>0\)).
\(\square\)

\paragraph{Consequence for Theorem~\ref{thm:core}.}
The compensator \(\tfrac12(f+Jf)\) ``produces the Hawking photon at \(\beta=0\)''.  With \(J=J_B\) it produces the value
\(0\) of the matter--tachyon coordinate \(\operatorname{Re}s=E^2-B^2\), not a null field.  The output keeps
\(E\cdot B\neq0\) and lies on one side of the photon, which is the endpoint of its image.  The statement holds for the real
part of \(s\) only.

\paragraph{Theoremlet v4.9.0-r02.5.}
\(J_B\) reverses the sense of the duality circle, \(J_BR_\alpha=R_{-\alpha}J_B\) with \(R_\alpha z=e^{i\alpha}z\).  Hence
\[
R_\alpha^{-1}J_BR_\alpha=R_{-2\alpha}J_B,
\qquad
R_{\pi/2}^{-1}J_BR_{\pi/2}=-J_B,
\qquad
R_{\pi/4}^{-1}J_BR_{\pi/4}=T .
\]
Passing from matter to the tachyon, reflecting there, and passing back gives the direct reflection followed by the
round-trip Umkehrwalze \(C\).

\paragraph{Proof sketch.}
\(e^{-i\alpha}\,i\,\overline{e^{i\alpha}z}=ie^{-2i\alpha}\bar z\).  Checked numerically.  \(\square\)

\paragraph{Reading.}
This is the Enigma pattern of the computational interpretation surface: a reflector between a forward and a backward
passage.  The two Umkehrwalzen of this layer are related by it.  The central one, \(C\), is the holonomy of a passage
pair.  The involutive one, \(J_B\), is the reflector, and conjugating it by the matter--tachyon passage multiplies it by
\(C\).

\subsection*{v4.9.0-r02 what the OP5 layer contributes}
\label{subsec:v490-r02-op5}

Inserting \(J=J_B\) into the leakage diagnostics \(E_{\rm comm}=[\mathcal M,J]\) and
\(E_{\rm leak}=(\mathrm{id}-C_{\HC})\mathcal MC_{\HC}\) of Definition~\ref{def:working-shared-matrix-leakage} gives the
following, with the transports of Reading~B in place of \(\mathcal M\) (relative norms on a random field):
\begin{center}
\small
\begin{tabular}{@{}lrr@{}}
\toprule
transport \(\mathcal M\) & \(|E_{\rm comm}z|/|z|\) & \(|E_{\rm leak}z|/|C_{\HC}z|\)\\
\midrule
real rotation & \(0\) & \(0\)\\
boost, \(\eta=0.6174\) & \(0.985\) & \(0.526\)\\
duality rotation, \(\alpha=0.4\) & \(0.779\) & \(0.389\)\\
\bottomrule
\end{tabular}
\end{center}
By Proposition~\ref{prop:working-stable-channel-criterion} the admissible channel is stable under the compact part of
the transport only.  Boosts leak because \(J_B\) singles out a frame.  Duality rotations leak because \(J_B\) reverses
their sense (Theoremlet r02.5).  Boosts preserve the side of the field as a whole, so this leakage is into the
\(J_B\)-odd component and does not change side.

\paragraph{Correspondence, offered and not adopted.}
The local OP5 closure theorem works under the identification of the OP5 near-kernel branch with the
\emph{photon-frame HC-cone/reflector route}.  Its halo invariance off tunnel candidates asks for a turning-point
readout before a through-collapse exit.  In the language of this layer a turning point is case (a) of Theoremlet r02.1,
a return on the same side with trivial holonomy, and a through-exit is case (b), a change of side with holonomy \(C\).
The reflector of that route would be antilinear by Theoremlet r02.3.  The radial and \(\HC\)-residual channels
\(\rho_r,\rho_h\), the kernel \(K\) and the tunnel selector have no counterpart in Reading~B yet.  The correspondence is
therefore recorded and not used.  It does not touch the guarded scope of the OP5 closure.

\subsection*{v4.9.0-r02 register entries}
\label{subsec:v490-r02-register}

\begin{center}
\small
\begin{tabular}{@{}>{\raggedright\arraybackslash}p{0.34\textwidth}>{\raggedright\arraybackslash}p{0.58\textwidth}@{}}
\toprule
statement & status after v4.9.0-r02\\
\midrule
Umkehrwalze of the round trip & holonomy of a side-changing passage pair, equal to \(C=PT\) on field strengths (r02.1, r02.2)\\
involutive \(J\) exchanging matter and tachyon & must be antilinear (contains \(T\)); representative \(J_B\): \(E\leftrightarrow B\) (r02.3); frame-dependent\\
Thm.~\ref{thm:core}, photon from \(\tfrac12(f+Jf)\) & with \(J_B\): \(\operatorname{Re}s=0\), not a null field; holds for the real part of \(s\) only (r02.4)\\
Def.~\ref{def:working-shared-matrix-leakage} with \(J_B\) & \(E_{\rm comm}=E_{\rm leak}=0\) for real rotations only (r02)\\
OP5 turning point / through-exit & correspondence to cases (a)/(b) of r02.1 recorded, not adopted\\
\bottomrule
\end{tabular}
\end{center}

\subsection*{v4.9.0-r02 ledger and guard update}
\label{subsec:v490-r02-ledger-guard}

\[
R_{\rm uw}^{v490\text{-}r02}=(0,0,1)
\]
for (round-trip Umkehrwalze identified with \(C\), involutive Umkehrwalze characterised as antilinear,
identification with the OP5 halo/tunnel structure open).

The layer spends no guard:
\[
R_{\rm glob}^{(v490\text{-}r02,+)}=(0,1,1,1,1),
\qquad
\Gamma_{\rm guard}^{v490\text{-}r02}=(1,1,1,1,1),
\qquad
\beta_4=1/256\ \text{unspent}.
\]
It does not imply exact finite \(C5\) zero closure, exact all-readout terminality,
\(C6/T2\), a global photon-clock/RH theorem, or global all-readout closure.

